\uuid{zfYx}
\titre{Valeur absolue}
\auteur{bodin ; hanani ; mzari}
\organisation{exo7}
\licence{CC BY-NC-SA 4.0}
\datecreate{2020-01-20}
\niveau{L1}
\module{Analyse}
\chapitre{Propriétés de R}
\sousChapitre{Propriétés des nombres réels}
\difficulte{2}
\theme{Propriétés de nombres réels, valeur absolue}
\exo7qcmid{lille-1.reels.010}

\contenu{
\begin{question}
\qclassification{120.03}
\qtags{motcle=valeur absolue}
Pour $x\in\Rr$, on définit $f(x)= x - |x|$.
Quelles sont les assertions vraies ?
\begin{answers}
    \bad{$\forall x \in \Rr \quad f(x)\ge0$}

    \good{$\forall x \in \Rr \quad f(x)\le0$}

    \good{$\forall x >0 \quad f(x) =  0$}

    \bad{$\forall x < 0 \quad f(x) =  -2x$}

\end{answers}
\begin{explanations}
Si $x \ge 0$, alors $f(x)=0$. Si $x<0$ alors $-|x|=x$ et donc $f(x)=2x<0$.
\end{explanations}
\end{question}
}
